%!TEX root = ../template.tex
%%%%%%%%%%%%%%%%%%%%%%%%%%%%%%%%%%%%%%%%%%%%%%%%%%%%%%%%%%%%%%%%%%%%
%% 5-Conlcusions.tex
%% NOVA thesis document file
%%
%% Chapter with the template manual
%%%%%%%%%%%%%%%%%%%%%%%%%%%%%%%%%%%%%%%%%%%%%%%%%%%%%%%%%%%%%%%%%%%%

\typeout{NT FILE 5-Conlcusions.tex}%

\chapter{Timeline and Conclusions}
\label{cha:Timeline and Conclusions}

\glsresetall

\section{Timeline} \label{sec.5:Timeline}
The timeline for the thesis is shown in Figure~\ref{fig.5:Monthly view}. On a monthly basis, one week per month is reserved for writing the thesis. If the writing task requires less than one week, the remaining time can be allocated to subsequent tasks.  

From a broader project perspective, shown in Figure~\ref{fig.5:Project view}, there are three periods during which the primary focus is on writing, either for the preparation of papers or for the final thesis writing phase. As before, if less time is required for writing, the schedule can be adjusted flexibly depending on the needs of the project at that time.

\begin{figure}[H]
    \begin{center}
    \begin{ganttchart}[y unit title=0.4cm,
    y unit chart=0.5cm,
    vgrid,hgrid, 
    title label anchor/.style={below=-1.6ex},
    title left shift=.05,
    title right shift=-.05,
    title height=1,
    progress label text={},
    bar height=0.7,
    group right shift=0,
    group top shift=.6,
    group height=.3]{1}{20}
    %labels
    \gantttitle{Week 1}{5}
    \gantttitle{Week 2}{5}
    \gantttitle{Week 3}{5}
    \gantttitle{Week 4}{5}\\
    %tasks
    \ganttbar{task at hand}{1}{15} \\
    \ganttbar{write thesis}{16}{20}
    
    %relations 
    \ganttlink{elem0}{elem1} 
    \end{ganttchart}
    \end{center}
    \caption{Monthly view}
    \label{fig.5:Monthly view}
\end{figure}

\begin{figure}[H]
    \hspace{-2.8cm}
    \begin{ganttchart}[y unit title=0.4cm,
    y unit chart=0.5cm,
    vgrid,hgrid, 
    title label anchor/.style={below=-1.6ex},
    title left shift=.05,
    title right shift=-.05,
    title height=1,
    progress label text={},
    bar height=0.7,
    group right shift=0,
    group top shift=.6,
    group height=.3]{1}{32}
    %labels
    \gantttitle{Feb.}{4}
    \gantttitle{Mar.}{4}
    \gantttitle{Apr.}{4}
    \gantttitle{May}{4}
    \gantttitle{June}{4}
    \gantttitle{July}{4}
    \gantttitle{Aug.}{4}
    \gantttitle{Sept.}{4}\\
    %tasks
    \ganttbar{agent to nearest point}{1}{4} \\
    \ganttbar{Paper 1}{5}{8} \\
    \ganttbar{Consensus}{9}{12} \\
    \ganttbar{Poisson disk sampling}{13}{16} \\
    \ganttbar{Other SO(2) Manifolds}{17}{20} \\
    \ganttbar{Paper 2}{21}{24} \\
    \ganttbar{IRL testing}{25}{28} \\
    \ganttbar{write thesis}{29}{32}
    
    %relations 
    \ganttlink{elem0}{elem1} 
    \ganttlink{elem0}{elem2} 
    \ganttlink{elem2}{elem3} 
    \ganttlink{elem3}{elem4} 
    \ganttlink{elem4}{elem5} 
    \ganttlink{elem4}{elem6} 
    \ganttlink{elem6}{elem7} 
    \end{ganttchart}
    \caption{Project view}
    \label{fig.5:Project view}
\end{figure}


\section{Conclusion}

This section has presented a comprehensive overview of circular formation control for multi-agent systems. We reviewed fundamental problem formulations, covering various agent dynamics, formation objectives, and classification schemes. Key control methodologies including cyclic pursuit, leader-follower, behavior-based, artificial potential, and virtual structure approaches were examined in detail.

The section highlighted the importance of measurement types—position, distance, bearing, and range—in shaping control strategies. Theoretical foundations including stability analysis, convergence proofs, and collision avoidance were discussed. Practical considerations such as input constraints, communication limitations, and robustness to disturbances were addressed.

Recent advances in encirclement-guaranteed pursuit strategies using robust model predictive control demonstrate the evolution of the field toward more sophisticated and application-oriented solutions. The ability to handle arbitrary phase distributions, moving targets, and complex constraints significantly expands the applicability of circular formation control.

Looking forward, several exciting research directions remain open. Three-dimensional formations, heterogeneous agent teams, obstacle-rich environments, and security concerns present significant challenges. The integration of learning-based methods with classical control theory offers promising avenues for handling uncertainty and adaptation.

As autonomous systems become increasingly prevalent in aerospace, marine, and ground applications, circular formation control will continue to play a vital role in enabling coordinated multi-agent operations. The theoretical foundations and practical solutions reviewed in this section provide a solid basis for future innovations in this important area.